\section{Discussion and limitations}
\label{sec:discussion}

\paragraph{What changes when the program is executed.} The largest held-out improvement to the optimizer came from no new geometry at all: delivering the program that actually builds valid recovered most of what the search had found. Under the same execution-based judge, a frontier model with inspection, plain code and execution feedback reached the optimizer's level on reserved held-out parts, at a few minutes and cents per part. The optimizer's stages first looked useless as agent tools; profiling located the cost in redundant computation, and with a revised package (exact rewrites plus a per-check time cap) those stages gave the strongest agent a small, consistent edge over plain code. Measuring tools by their cost as well as their output therefore changed the conclusion about them; which part of the package matters was not isolated.

\paragraph{Limitations.} (i) The optimizer study's measurement was amended after outcomes were visible (one terminal recovery evaluation for setup and resource-limit replay failures); the original analysis is reported beside it, and it changes the dx and v2 values substantially. (ii) The agent studies use 40 parts per cohort with three episodes each, so many intervals are several hundredths of IoU wide, and the Study 2 comparisons are exploratory. (iii) One benchmark: hard-tier CADBench parts from two splits. (iv) The endpoint is volumetric IoU; surface agreement and an independent re-judging are pending. (v) The code toolset includes target inspection and execution feedback, so it measures an agent in this harness, not a language model alone. (vi) Closed-model results are tied to specific model versions, settings (32k tokens per turn, full history) and an API; the tools-v3 package bundles exact rewrites with a behavioural check cap, and its output-equivalence was checked on development inputs only. (vii) Infrastructure failures (faulty nodes, API limits) required reruns, listed in \cref{app:log}.
